
\documentclass[10pt]{article}
\usepackage[utf8]{inputenc}
\usepackage[T1]{fontenc}
\usepackage{amsmath, amssymb, xcolor, geometry, enumitem, multicol, tcolorbox}

\definecolor{mainblue}{RGB}{0, 102, 204}
\geometry{a4paper, margin=1.2cm, top=1cm, bottom=3cm, footskip=1.2cm}

\begin{document}
\enlargethispage{2.5cm}

% Modern Header
\begin{tcolorbox}[colback=mainblue!5, colframe=mainblue, sharp corners, boxrule=0.5pt]
    \small \textbf{Unit:} Unit 0: Foundations of Algebra \hfill \textbf{Name:} \rule{3cm}{0.4pt} \\
    \small \textbf{Topic:} GCF & LCM \hfill \textbf{Date:} \rule{2cm}{0.4pt} \\
    \centering \textbf{\large Foundation (Jalapeno)}
\end{tcolorbox}

\vspace{0.2cm}

% MCQ Section
\begin{multicols}{2}
\begin{enumerate}[label=\textbf{Q\arabic*.}, leftmargin=*, itemsep=6pt, parsep=0pt]

  \item \small Find the GCF of $12$ and $18$ using factors.
  \begin{enumerate}[label=(\Alph*), leftmargin=0.6cm, nosep, itemsep=2pt]
    \item $1$
    \item $2$
    \item $3$
    \item $6$
    \item $9$
  \end{enumerate}

  \item \small What is the LCM of $4$ and $6$ using multiples?
  \begin{enumerate}[label=(\Alph*), leftmargin=0.6cm, nosep, itemsep=2pt]
    \item $4$
    \item $6$
    \item $8$
    \item $12$
    \item $24$
  \end{enumerate}

  \item \small Find the prime factorization of $45$ to determine its GCF with $15$.
  \begin{enumerate}[label=(\Alph*), leftmargin=0.6cm, nosep, itemsep=2pt]
    \item $3 \cdot 5$
    \item $3 \cdot 3 \cdot 5$
    \item $3^2 \cdot 5$
    \item $3 \cdot 5^2$
    \item $3^2 \cdot 5^2$
  \end{enumerate}

  \item \small A business has stocks priced at $12$ and $18$ per share. What is the LCM of these prices?
  \begin{enumerate}[label=(\Alph*), leftmargin=0.6cm, nosep, itemsep=2pt]
    \item $12$
    \item $18$
    \item $36$
    \item $54$
    \item $108$
  \end{enumerate}

  \item \small Determine the GCF of $24$ and $30$ using prime factorization.
  \begin{enumerate}[label=(\Alph*), leftmargin=0.6cm, nosep, itemsep=2pt]
    \item $1$
    \item $2$
    \item $3$
    \item $6$
    \item $12$
  \end{enumerate}

  \item \small A savings account earns $3\%$ and $6\%$ interest on two separate investments. Find the GCF of these interest rates.
  \begin{enumerate}[label=(\Alph*), leftmargin=0.6cm, nosep, itemsep=2pt]
    \item $1\%$
    \item $2\%$
    \item $3\%$
    \item $6\%$
    \item $9\%$
  \end{enumerate}

  \item \small What is the LCM of $8$ and $10$ using prime factorization?
  \begin{enumerate}[label=(\Alph*), leftmargin=0.6cm, nosep, itemsep=2pt]
    \item $8$
    \item $10$
    \item $20$
    \item $40$
    \item $80$
  \end{enumerate}

  \item \small Determine the LCM of $9$ and $12$ to calculate the total amount of tax owed.
  \begin{enumerate}[label=(\Alph*), leftmargin=0.6cm, nosep, itemsep=2pt]
    \item $9$
    \item $12$
    \item $18$
    \item $27$
    \item $36$
  \end{enumerate}

\end{enumerate}
\end{multicols}

\vspace{-0.2cm}
\hrule
\vspace{0.2cm}
\textbf{\small Open Ended Questions (FRQ)}
\begin{enumerate}[label=\textbf{Q\arabic*.}, start=9, itemsep=5pt, topsep=0pt]

  \item \small A company has two investments, one with a $5\%$ annual return and another with a $10\%$ annual return. Find the LCM of these interest rates and explain how it can be used to calculate the total return on investment over a $2$-year period. \vspace{2cm}

  \item \small In a stock market, two stocks have prices of $15$ and $25$ per share. Find the GCF and LCM of these prices and discuss how they can be used to make informed investment decisions. Show your calculations and provide a brief explanation. \vspace{2cm}

\end{enumerate}

\vfill

% Chef's Tips Footer
\begin{tcolorbox}[colback=blue!5, colframe=mainblue!50, title=\small \textbf{Chef's Tips}, fonttitle=\bfseries, bottom=1mm]
\begin{itemize}[noitemsep, topsep=2pt, leftmargin=0.5cm]
  \item \scriptsize Tip 1: Encourage students to use prime factorization to find GCF and LCM.\item \scriptsize Tip 2: Emphasize the importance of understanding the concept of GCF and LCM in real-world applications.\item \scriptsize Tip 3: Provide feedback on the free-response questions to help students improve their problem-solving skills.
\end{itemize}
\end{tcolorbox}

\end{document}
  